\section{总结与延伸}

本讲把 frontier scaling 从“模型越大越好”扩展为一套 research、systems、safety 与 product operating system：

\begin{enumerate}
  \item \textbf{Scaling law 是预测合同}：先用多尺度 run 做 fit 和 residual，再批准 frontier commitment；
  \item \textbf{Compute multiplier 需要归因}：model、data、optimization 与 systems 改动要在共同成本基准下验证；
  \item \textbf{Rare failure 必须被设计进去}：failure domain、checkpoint、replay 和 recovery tax 决定 useful progress；
  \item \textbf{Training 需要四平面可观测性}：model、optimizer、data、system telemetry 共同解释 anomaly；
  \item \textbf{全球值班需要 decision rights}：handoff packet、owner 与 stop criteria 比“有人在线”更重要；
  \item \textbf{Post-training 优化 proxy}：RLHF/CAI/RLAIF 扩大行为监督，但都需要独立 evaluation 与治理；
  \item \textbf{Evaluation 取决于 elicitation}：task、tool、effort、environment、judge 和 uncertainty 决定分数含义；
  \item \textbf{Safeguard 随 capability 升级}：RSP 版本会演化，先评估、满足 required controls、再训练/部署的原则不变；
  \item \textbf{Interpretability 是补充证据}：feature/circuit 需要 causal validation，不能单独证明安全；
  \item \textbf{API 是外部 contract}：chat 可快速试验，API 必须 version、migrate、deprecate 与 support。
\end{enumerate}

\begin{importantbox}{Frontier Model Program 作业}
设计一个 8 周训练与发布计划：给出三档 pilot scale、loss fit 与 uncertainty；列出 workers/network/storage/scheduler failure budget；定义 checkpoint/replay 和 follow-the-sun handoff；画出 pre-training、RLHF/RLAIF、capability/safety evaluation、safeguard gate、chat canary 和 API release；最后给出 model deprecation、shadow traffic 与 rollback。每个 gate 必须有 owner、evidence 和 stop condition。
\end{importantbox}

\subsection{拓展阅读}

{\small
\begin{itemize}[nosep,leftmargin=*]
  \item \href{https://arxiv.org/abs/2001.08361}{\textbf{Scaling Laws for Neural Language Models}}：经验缩放与预测方法。
  \item \href{https://arxiv.org/abs/2203.02155}{\textbf{InstructGPT}}：RLHF pipeline 与限制。
  \item \href{https://www.anthropic.com/research/constitutional-ai-harmlessness-from-ai-feedback}{\textbf{Constitutional AI}}：principles、critique/revision 与 RLAIF。
  \item \href{https://www.anthropic.com/research/evaluating-ai-systems}{\textbf{Challenges in evaluating AI systems}}：评估设计与现实有效性。
  \item \href{https://www.anthropic.com/responsible-scaling-policy/rsp-v3-0}{\textbf{Anthropic RSP v3.0}}：2026 年当前 capability-threshold framework。
  \item \href{https://www.anthropic.com/research/towards-monosemanticity-decomposing-language-models-with-dictionary-learning}{\textbf{Towards Monosemanticity}}：dictionary learning 与 feature。
  \item \href{https://docs.anthropic.com/en/docs/about-claude/model-deprecations}{\textbf{Claude model lifecycle}}：active、legacy、deprecated 与 retired。
\end{itemize}
}

\end{document}
